\documentclass[11pt]{article}

\usepackage[margin=0.75in]{geometry}
\usepackage[T1]{fontenc}
\usepackage{lmodern}
\usepackage{booktabs}
\usepackage{tabularx}
\usepackage{array}
\usepackage{amsmath}
\usepackage{xcolor}
\usepackage[hidelinks]{hyperref}

\newcommand{\dflash}{DeepSeek V4 Flash}
\newcommand{\dpro}{DeepSeek V4 Pro}
\newcommand{\dflashshort}{DS V4 Flash}
\newcommand{\dproshort}{DS V4 Pro}
\newcommand{\sol}{GPT-5.6 Sol}
\newcommand{\na}{--}
\newcolumntype{Y}{>{\raggedright\arraybackslash}X}

\title{Reproducing SkillOpt on VeruSAGE:\\
An Epoch-1 Comparison of Actor and Optimizer Configurations}
\author{}
\date{August 13, 2026}

\begin{document}
\maketitle

\begin{abstract}
This document summarizes our current reproduction of SkillOpt on VeruSAGE
proof-repair tasks. The reproduced system evolves one monolithic skill and
does not perform retrieval at inference time. We use a frozen split of 40
training problems, 20 selection problems, and 40 held-out test problems. A
first epoch evaluates the initial skill on the selection set, collects 40
training rollouts, updates the skill with an optimizer, and evaluates the
candidate on the same selection set. Later epochs reuse the score of the
current accepted skill and therefore require only 40 training rollouts and 20
selection-gate rollouts. The candidate is accepted only if its selection
accuracy strictly improves. With \dflash{} as the actor, both the
original Flash optimizer and the stronger \sol{} optimizer produced candidates
that decreased selection accuracy from 6/20 to 4/20. A \dpro{} optimizer
candidate was rejected by an evidence audit before live evaluation. In the
latest complete epoch, using \dpro{} as the actor and \sol{} as the optimizer
improved selection accuracy from 16/20 to 17/20, giving the first accepted
candidate in this reproduction. This is an in-loop selection result; the
40-problem held-out test set has not been evaluated.
\end{abstract}

\section{Reproduction Scope}

The current SkillOpt reproduction maintains a single textual skill. The whole
skill is supplied to the actor for every rollout, and the optimizer revises
that skill using actor execution traces. The current version does not contain
a runtime retrieval component and therefore does not select task-specific
skill clauses or memories. The experiment studied here is consequently a
reproduction of iterative global-skill optimization, rather than a retrieval-
augmented extension.

SkillOpt separates model use into two roles:

\begin{itemize}
    \item \textbf{Actor:} receives the current skill and executes a VeruSAGE
    proof-repair task. Its trajectory, verifier outcome, and final candidate
    are recorded.
    \item \textbf{Optimizer:} analyzes the training trajectories, reflects on
    successes and failures, merges proposed changes, ranks the changes, and
    produces the next candidate skill.
\end{itemize}

The standard SkillOpt configuration uses GPT-5.5 for both roles. Our
experiments vary the actor and optimizer independently to determine whether
the negative Flash results are associated with the optimizer, the actor
trajectories, or both.

\section{Data Split and Epoch Workflow}

All reported experiments use the same frozen Anvil/IronKV split. Its SHA-256
identifier is
\begin{center}
\small
\texttt{53059264e5d0458e1fc50a3c1786cbea}\\
\texttt{c6c671aedf56dd71fb32843b24d2c553}.
\end{center}
The split contains 40 training problems, 20 selection problems, and 40
held-out test problems. The selection set is used only for in-loop skill
evaluation. Training rollouts are the only task trajectories supplied to the
optimizer. The held-out set is not used in the experiments summarized here.

For the first epoch of one continuous SkillOpt run, the actor workload is

\[
N_{\mathrm{actor}}
= N_{S_0} + N_{\mathrm{train}} + N_{\mathrm{gate}}
= 20 + 40 + 20
= 80 \text{ task rollouts}.
\]

The initial $S_0$ selection evaluation is a one-time initialization cost. Once
an accepted skill and its selection score have been established, each later
epoch requires only

\[
N_{\mathrm{later\ epoch}}
= N_{\mathrm{train}} + N_{\mathrm{gate}}
= 40 + 20
= 60 \text{ task rollouts}.
\]

Consequently, a continuous run of $E$ epochs requires

\[
N_{\mathrm{actor}}(E) = 20 + 60E,
\]

excluding any one-time held-out test evaluation. For example, two epochs use
140 actor rollouts and four epochs use 260. An epoch would return to 80 only
if it were launched as a separate experiment that deliberately reruns the
20-problem baseline, rather than continuing from the stored accepted score.

The workflow is shown in Table~\ref{tab:workflow}. Here, $S_0$ is the initial
accepted skill, and $S_1$ is the candidate produced by the optimizer.

\begin{table}[h]
\centering
\small
\begin{tabularx}{\textwidth}{@{}clcY@{}}
\toprule
Step & Stage & Actor rollouts & Purpose \\
\midrule
1 & Initial selection evaluation & 20 & Measure $S_0$ on the fixed selection set. \\
2 & Training rollout & 40 & Generate success and failure trajectories using $S_0$. \\
3 & Skill optimization & 0 & Reflect, merge, rank, and construct candidate $S_1$. \\
4 & Candidate selection gate & 20 & Measure $S_1$ on the same selection problems. \\
5 & Gate decision & 0 & Accept $S_1$ only if its strict solved rate exceeds $S_0$. \\
\bottomrule
\end{tabularx}
\caption{One-epoch SkillOpt workflow used in the reproduction.}
\label{tab:workflow}
\end{table}

For an optimizer replay, stored training trajectories can be reused. Such a
replay requires no new training rollouts; only the candidate's 20-problem gate
is rerun. We distinguish these replays from complete epochs throughout the
results.

\section{Experimental Configurations}

Four actor--optimizer configurations have been investigated:

\begin{enumerate}
    \item \textbf{Flash/Flash (F/F):} a complete epoch using \dflash{} as both
    actor and optimizer.
    \item \textbf{Flash/Pro (F/P):} offline reanalysis of the stored Flash
    training trajectories using \dpro{} as the optimizer. The final candidate
    was rejected before the live gate because deterministic and manual audits
    identified incorrect evidence attribution.
    \item \textbf{Flash/Sol (F/S):} native SkillOpt reflection, merge, and
    ranking over the same stored Flash trajectories using local Codex
    \sol{} as the optimizer, followed by a fresh 20-problem Flash actor gate.
    \item \textbf{Pro/Sol (P/S):} a complete epoch using \dpro{} as the actor
    through the Codex CLI native Responses harness and local Codex \sol{} as
    the optimizer.
\end{enumerate}

The initial skill is the same 838-byte artifact in all configurations. The
Flash replay experiments inherit the $S_0$ and training measurements from the
complete F/F epoch instead of regenerating them.

\section{Primary Results}

Table~\ref{tab:main-results} gives the raw solved counts. ``Gain/Loss'' is a
paired comparison on the fixed selection tasks: gain is fail-to-pass and loss
is pass-to-fail. The F/P row deliberately contains no candidate score because
the candidate never entered a live selection gate. Reporting it as 4/20 would
incorrectly convert a pre-gate audit rejection into an evaluated result.

\begin{table*}[t]
\centering
\small
\setlength{\tabcolsep}{4pt}
\begin{tabularx}{\textwidth}{@{}llYccccccl@{}}
\toprule
ID & Actor & Optimizer & $S_0$ & Train & Candidate & $\Delta$ & Gain/Loss & Skill bytes & Decision \\
\midrule
F/F & \dflashshort{} & \dflash{}
    & 6/20 & 8/40 & 4/20 & $-2$ & 0/2 & 10,322 & Reject \\
F/P & \dflashshort{} & \dpro{}
    & $6/20^{\dagger}$ & $8/40^{\dagger}$ & \na{} & \na{} & \na{} & 1,646 & Pre-gate reject \\
F/S & \dflashshort{} & \sol{}
    & $6/20^{\dagger}$ & $8/40^{\dagger}$ & 4/20 & $-2$ & 0/2 & 3,490 & Reject \\
P/S & \dproshort{} & \sol{}
    & 16/20 & 35/40 & \textbf{17/20} & $\mathbf{+1}$ & \textbf{1/0} & 2,932 & \textbf{Accept} \\
\bottomrule
\end{tabularx}
\vspace{2pt}
\begin{minipage}{\textwidth}\footnotesize
\emph{Notes.} $\Delta$ is the change in solved problems on the 20-problem
selection set. Gain/Loss reports fail-to-pass and pass-to-fail transitions,
respectively. $\dagger$ indicates that the Flash $S_0$ result and 40 training
trajectories were reused from the F/F epoch rather than rerun. No row includes
held-out test evaluation.
\end{minipage}
\caption{Selection-gate results for the SkillOpt reproduction.}
\label{tab:main-results}
\end{table*}

The Flash actor produced the same negative gate result with two different
evaluated optimizers. The F/F candidate introduced zero new solves and lost
two baseline solves. Replacing the optimizer with \sol{} produced a shorter,
audit-clean candidate, but the paired selection outcome remained exactly zero
gains and two losses. The Pro optimizer analysis cannot be included in this
numerical comparison because its candidate was stopped before target
evaluation.

The P/S epoch produced the first accepted update. The initial skill solved
16/20 selection tasks, the training rollout solved 35/40 tasks, and the
candidate solved 17/20 selection tasks. The paired result was one gain, zero
losses, 16 retained successes, and three retained failures. SkillOpt therefore
accepted the candidate because 17/20 strictly exceeded 16/20.

\section{Usage and Cost}

Table~\ref{tab:cost} separates newly executed actor work from optimizer work.
This distinction is important for replay experiments: F/P ran no actor tasks,
and F/S reused all 40 training trajectories and paid only for a new candidate
gate. The P/S actor ledger includes the full initial evaluation, training
rollout, and candidate gate.

\begin{table*}[t]
\centering
\scriptsize
\setlength{\tabcolsep}{4pt}
\begin{tabularx}{\textwidth}{@{}YYrrYYYr@{}}
\toprule
Actor & Optimizer & Tasks & Actor req. & Input (cached) & Output & Actor cost & Optimizer cost \\
\midrule
\dflashshort{} & \dflash{} & 80 & 4,184 & 35.528M (27.766M) & 14.403M & \$5.1971 & \$0.0351 \\
\dflashshort{} & \dpro{} & 0 & 0 & \na{} & \na{} & \$0 & \$0.0581 \\
\dflashshort{} & \sol{} & 20 & 1,627 & 12.925M (9.559M) & 5.447M & \$2.0230 & \$0 local \\
\dproshort{} & \sol{} & 80 & 3,559 & 197.107M (194.614M) & 2.986M & \$4.3879 & \$0 local \\
\bottomrule
\end{tabularx}
\vspace{2pt}
\begin{minipage}{\textwidth}\footnotesize
\emph{Notes.} Parentheses in ``Input (cached)'' report cache-hit input tokens.
The Flash/Pro row is optimizer-only reanalysis of stored Flash trajectories,
so it has no new actor tasks. Optimizer costs use model-specific prices
recorded at experiment time. Flash/Flash uses the Flash rates:
\$0.0028/M cache-hit input, \$0.14/M cache-miss input, and \$0.28/M output.
Flash/Pro uses the Pro rates: \$0.003625/M cache-hit input, \$0.435/M cache-miss
input, and \$0.87/M output. Because the optimizer logs do not retain a prompt
cache split, their DeepSeek prompt tokens are conservatively treated as cache
misses. The displayed Flash/Pro value is the final Pro-v2 analysis; including the
earlier rejected Pro-v1 analysis gives \$0.1126. F/S and P/S used local Codex
quota for the \sol{} optimizer and therefore had zero metered optimizer cash
cost. At the recorded Sol API rates and treating all optimizer input as
uncached, their API-equivalent costs would be approximately \$1.5671 and
\$5.7069, respectively, before any long-context uplift. P/S actor cost excludes
Pro-native preflights and a stopped wrong-reasoning startup. Including those
bring-up calls yields approximately \$4.7360 for that Pro launch.
\end{minipage}
\caption{Measured actor usage and optimizer usage. Dollar values are estimated
from the provider usage ledgers.}
\label{tab:cost}
\end{table*}

The Pro/Sol epoch used substantially more input tokens than Flash/Flash, but
98.7\% of its input was served from the DeepSeek prompt cache and it produced
only 2.986M output tokens. By comparison, Flash/Flash had a 78.2\% input cache
rate and produced 14.403M output tokens. Output alone cost approximately
\$4.0327 in Flash/Flash, versus \$2.5980 in Pro/Sol. The cache-hit, cache-miss,
and output contributions were therefore \$0.0777, \$1.0866, and \$4.0327 for
Flash/Flash, compared with \$0.7055, \$1.0844, and \$2.5980 for Pro/Sol. This
token-shape difference outweighs Pro's higher unit prices, yielding formal
actor costs of \$5.1971 and \$4.3879, respectively. This is a measured harness
outcome, not evidence that Pro is intrinsically cheaper. The local \sol{}
optimizer incurred no metered API cash cost in either Sol experiment.

\section{Interpretation}

\paragraph{Optimizer strength alone did not rescue the Flash actor.}
The evaluated F/F and F/S candidates both changed selection performance from
6/20 to 4/20, with identical paired transitions of zero gains and two losses.
The stronger Sol optimizer improved skill compactness and passed the skill
contract audit, but this did not translate into a better Flash actor gate.

\paragraph{The Pro actor generated more useful optimization evidence.}
In the P/S epoch, the Pro actor solved 35/40 training tasks and the optimizer
produced an accepted candidate with a paired $+1/0$ selection change. This is
consistent with the hypothesis that higher-quality actor trajectories can
provide more useful evidence for skill optimization. It is not, by itself, a
causal isolation of actor capability because the Pro experiment also used a
different, Codex-native target harness and a higher initial $S_0$ score.

\paragraph{The gate prevented harmful updates.}
Both evaluated Flash candidates were rejected and the accepted skill remained
$S_0$. The Pro candidate was accepted only after a strict improvement on the
same fixed selection tasks. This demonstrates the intended mechanical role of
the SkillOpt gate even though it does not establish population-level gains.

\paragraph{The Pro-optimizer result is qualitative, not a gate score.}
The F/P candidate was compact, but its evidence audit found incorrect causal
attribution of a verifier failure. Stopping it before the live gate was an
integrity decision. It should not be grouped numerically with the two 4/20
Flash gate results.

\section{Validity and Limitations}

The evidence has four boundaries. First, every selection comparison contains
only 20 tasks, so one task corresponds to five percentage points and sampling
uncertainty is large. Second, the experiments are single runs without an A/A
estimate of actor stochasticity. Third, the old Flash and new Pro experiments
share the frozen split and gate semantics but not an identical target harness;
the latest positive result cannot be attributed solely to changing the actor
model. Fourth, the selection set is reused for in-loop acceptance and is not a
held-out generalization set. No claim of improved held-out solved rate or token
efficiency is made here.

All 80 final task results in the P/S epoch completed with 77 full V2 traces,
three timeout-preserved V1 traces, zero invalid V0 results, zero provider
errors, and unchanged source inputs. The 40 held-out test tasks were not run in
any experiment summarized by the primary table.

\section{Conclusion}

The reproduction confirms the core mechanics of single-skill SkillOpt on a
40/20/40 VeruSAGE split. With \dflash{} as actor, changing the optimizer from
Flash to \sol{} did not change the negative selection outcome: both evaluated
candidates moved from 6/20 to 4/20 and were rejected. The \dpro{} optimizer
candidate was rejected before live evaluation and therefore has no selection
score. The latest complete P/S epoch, using \dpro{} as actor and \sol{} as
optimizer, moved from 16/20 to 17/20 with one paired gain and no paired loss.
This candidate passed the first gate and became the first accepted evolved
skill in the current reproduction. The result is evidence of a successful
in-loop update, not held-out effectiveness evidence.

\end{document}
